% Build the layer, Chapter 2, second build: who is presenting a card.
% Inserted in \S\ref{sec:adversary}, after the paragraph on signing each
% request with the holder's key, which this section builds.
%
% No code or output in this file is typed by hand. \pdexcerpt, \pdrun and
% \pdrunlines print files under excerpts/, which the kernel's
% tools/excerpts.mjs cuts from its source and its recorded runs (see
% figures/pd-listing.tex). To change a listing, change the kernel and cut
% again. The prose names functions, never line numbers, so that a listing
% which moves in its file does not make a sentence wrong.
%
% \pdlistingroom{n} before a listing of n lines keeps it on one page. When a
% listing changes length, change the number with it.
\subsection{Build the layer: proof of possession}
\label{sec:anchor-build-holder}

\noindent\emph{Express lane: the gate of \S\ref{sec:anchor-build} admits
whoever presents valid cards, as their holder. Here the holder must sign
something new at every presentation. Two repairs that look correct are built
and broken first. The one that works does a check and a change in a single
step, as the cell of Chapter~0 did.}

The build of \S\ref{sec:anchor-build} stopped at a card that is a bearer
credential, and the paragraph above says what would fix that. This section
builds it, in \pdkernelname{} again and in the same four parts. Its
\emph{Naive} is where the first build finished.

\subsubsection{Naive: whoever presents the cards}

The verifier of Listing~\ref{alg:anchor-build-admit} checks the cards. It
checks nothing about whoever presents them. \pdkernelname{} keeps that gate
under the name \texttt{bearer}, so that it can be run.

\subsubsection{Break: a copied card}

Heron keeps his cards in a file, as an agent must keep them somewhere. Wren
reads the file.

\pdlistingroom{17}
\pdrun[caption={A copied card. Wren presents Heron's cards and is admitted
as Heron.},label={alg:anchor-holder-theft}]{theft-bearer}

Nothing in Listing~\ref{alg:anchor-holder-theft} is forged. Every check of
\S\ref{sec:anchor-build} passes, because the cards are Heron's own. The
gate was presented the right cards by the wrong party, and it had no way to
tell which party that was.

\subsubsection{A first repair: sign the request}

A card already names its holder's public key, since the next card in a
chain must be signed with it. So a holder can prove that it has the private
half of that key. It signs its request, and the gate verifies the signature
with the key the last card names. Listings~\ref{alg:anchor-holder-prove}
and~\ref{alg:anchor-holder-signed} are the two sides.

\texttt{proofText} builds the text that is signed, a list of four items:
the word \texttt{proof}, a challenge, the requested right, and the id of the
card the request is made under. The first repair has no challenge yet and
leaves it empty. The id is there for the reason \S\ref{sec:anchor-build}
gave for \texttt{par}. A holder may have two cards that name one key, so its
signature must name the card.

\pdlistingroom{24}
\pdexcerpt[caption={The holder's side. A proof is a signature made for one
presentation of its cards.},label={alg:anchor-holder-prove}]{ch02b-prove}

\pdlistingroom{18}
\pdexcerpt[caption={The first repair. The cards are checked as before, and
then the signature on the request.},label={alg:anchor-holder-signed}]{ch02b-signed}

Wren's copied cards are refused now. She can sign, but only with her own
key, and the last card names Heron's.

\subsubsection{Break again: the request is sent twice}

Wren does not need Heron's key. She needs one message that was signed with
it.

\pdlistingroom{16}
\pdrun[caption={A replay. What Heron sent is accepted a second
time.},label={alg:anchor-holder-replay}]{replay-signed}

The signature in Wren's message is Heron's. It covers the right she
requests and the card she presents. The gate cannot know that it has seen
this message before. A signature shows that a key was used. It does not show
when.

\subsubsection{A second repair: a challenge, looked up and then deleted}

A recorded message cannot contain the answer to a question that had not
been asked when it was recorded. So the gate asks one. Before each
presentation it issues a \emph{challenge}, sixteen random bytes, and the
holder signs the challenge together with the request. The gate keeps the
challenges it has issued in a table, and deletes one when its answer is
admitted. The protocols that introduced this called such a number a nonce, a
number used once~\pdcite{needham1978using}.

A gate serves many agents, and it does not make them all wait while it
verifies one signature. Listing~\ref{alg:anchor-holder-lookup} checks a
presentation in two steps, and other requests may run between them.

\pdlistingroom{14}
\pdexcerpt[caption={The second repair. The challenge is looked up in one
step and deleted in another.},label={alg:anchor-holder-lookup}]{ch02b-two-step}

The \texttt{yield} marks the pause, as it did in the workers of
\weld{prereq:sec:prereq-build}{yield}. Run one presentation after the other
and this gate refuses the replay: by the time Wren sends it, Heron's
challenge has been deleted.

\subsubsection{Break a third time: two presentations at once}

Wren sends her copy while the gate is still checking Heron's own message.
The gate takes two turns over each message, and the turns can fall in six
orders. \pdkernelname{} tries them all.

\pdlistingroom{15}
\pdrun[caption={Every order of the gate's turns over two presentations of
one proof.},label={alg:anchor-holder-showings}]{presentations-two-step}

In four orders both messages are looked up before either challenge is
deleted, and both are admitted. Four of six is the count of
Chapter~0's first break, and the cause is the same: a check and a change
with a pause between them
(\weld{prereq:sec:prereq-build}{lose an update}). The two safe orders are
the ones in which the gate finishes with one message before it starts on
the other. \emph{Now you try:} Exercise~\ref{ex:anchor-holder-orders} asks
why those two have three turns and not four.

\subsubsection{Fix: one decision}

The fix has two parts. The gate looks the challenge up and deletes it in
the same turn as the decision (Listing~\ref{alg:anchor-holder-gate}). A
function in the decision core, \texttt{held}, makes that decision after
\texttt{admit} has accepted the cards
(Listing~\ref{alg:anchor-holder-held}).

\pdlistingroom{27}
\pdexcerpt[caption={The gate. The lookup and the deletion are in one turn
with the decision.},label={alg:anchor-holder-gate}]{ch02b-gate}

\texttt{admit} has no \texttt{yield} in it and waits on nothing. From the
lookup to the deletion nothing else runs, so two presentations of one proof
cannot both find the challenge open. The core keeps no table, so the gate
passes it one row: the challenge that the proof names, if it is open.

\pdlistingroom{35}
\pdexcerpt[caption={The core's check of who is presenting a
chain.},label={alg:anchor-holder-held}]{ch02b-held}

A proof that names any other challenge is refused, and so is one whose
challenge was issued more than thirty seconds ago. Then the core computes
for itself the text that the proof must be a signature over
(Listing~\ref{alg:anchor-holder-text}). It takes the challenge from the
gate's row and the rest from the presentation it is deciding. From the proof
it takes only the signature. The pass that \texttt{held} returns names the
challenge that was answered.

\pdlistingroom{8}
\pdexcerpt[caption={What a proof is a signature
over.},label={alg:anchor-holder-text}]{ch02b-text}

The text begins with the word \texttt{proof}, and the signed text of a
card begins with the harbor's header. A gate, or someone impersonating it,
can send a holder anything at all as a challenge. Whatever it sends, what
the holder signs in answer cannot be presented as a card.

% A page may break here. See the note in anchor-build-the-verifier.tex.
\pagebreak[0]
\keyidea[A proof answers a question that had not been asked]{A copy of a
card is as good as the card, and a replay of a signature is as good as
the signature. What neither can contain is the answer to a question nobody
had asked when it was made. So the gate asks one, and accepts each answer
once.}

\subsubsection{Verify: the same runs, then every order}

The copied card first (Listing~\ref{alg:anchor-holder-theft-core}).

\pdlistingroom{13}
\pdrunlines[caption={The copied card, against the fix. Wren has the cards
and no key.},label={alg:anchor-holder-theft-core}]{theft-core}{1,13-24}{1 and 13\textendash 24}

Then the replay (Listing~\ref{alg:anchor-holder-replay-core}). Wren
tries it twice. The second time she requests a challenge of her own and
sends Heron's old signature as her answer to it.

\pdlistingroom{15}
\pdrunlines[caption={The replay, against the fix. Heron's signature was
made for another challenge.},label={alg:anchor-holder-replay-core}]{replay-core}{1,9-22}{1 and 9\textendash 22}

Then both messages at once (Listing~\ref{alg:anchor-holder-showings-core}).

\pdlistingroom{11}
\pdrun[caption={Every order, against the fix. The gate takes one turn over
a message.},label={alg:anchor-holder-showings-core}]{presentations-core}

There are two orders now, and the proof is admitted once in each. Look
at the second. Wren's copy arrives first, and it is the copy that is
admitted, while Heron's own message is refused. The gate cannot tell the
two apart, since both carry Heron's signature. What it guarantees is
narrower than ``only Heron'': his signature is accepted once, for the
request he signed. The note that closes this section returns to that.

Next, what the signed text has to contain
(Listing~\ref{alg:anchor-holder-proofs}).

\pdlistingroom{18}
\pdrun[caption={Which attacks succeed when the signed text leaves one part
out.},label={alg:anchor-holder-proofs}]{proofs}

The first row is the core's own gate. Each other row is a gate built like
it whose text leaves one part blank, and each admits one changed
presentation. Without the challenge, Heron's signature answers a challenge
that Wren requested herself, which is the column \texttt{replayed}. Without
the right, a request to read \texttt{shop/checkout.ts} can be forwarded as a
request to write it. The last row is the moved card of
\S\ref{sec:anchor-build} one level down: a request that Alice signed under
her narrow card is presented under her wide one. The core has no guard for
these that could be removed, because it never compares: it computes the
text itself. These gates are the control in its place.

Last, the guards that can be removed. Remove the deletion from the gate and
both presentations are admitted in both orders
(Listing~\ref{alg:anchor-holder-showings-mutant}).
Four more builds each remove one guard from the core: the requirement of a
proof, the check of its signature, the check that its challenge is open,
and the limit of thirty seconds. For each, the kernel names a test that
passes on the real core and fails on that build.

\pdlistingroom{11}
\pdrun[caption={The fix with the deletion removed from its
gate.},label={alg:anchor-holder-showings-mutant}]{presentations-keep-challenge}

The protocol is also three ProVerif models
(Listing~\ref{alg:anchor-holder-models}).

\pdlistingroom{13}
\pdrun[caption={ProVerif~2.05 on models of a holder answering a
challenge.},label={alg:anchor-holder-models}]{models-possession}

\texttt{Signed} marks the holder signing a proof, and \texttt{Accepted}
marks the gate accepting one. The first query of each model says that
whatever the gate accepts, the holder signed. The query that starts
\texttt{inj-event} says more: nothing is accepted more often than it was
signed. It holds for \texttt{answered-once.pv}, which is the fix. It fails
for the two repairs. In \texttt{no-challenge.pv} one signed request is
accepted twice, and in \texttt{challenge-left-open.pv} one answer is. In all
of them the attacker chooses what the holder signs and sees every signature.
\pdscope{The runs cover two presentations of one proof at one gate. The
models cover any number of presentations, with cryptography that cannot be
broken and no clock. The fix's model assumes that a gate accepts one answer
to a challenge, so the model cannot show that the code has no pause
between its lookup and its deletion. The run of every order shows that.}

\begin{pdboundary}
\setlength{\parskip}{4pt plus 1pt}
A proof is accepted once, for whoever delivers it first: that is the second
order of Listing~\ref{alg:anchor-holder-showings-core}. If Wren can
intercept Heron's message before the gate receives it, she is admitted
in his place, and the reply to his request goes to her. She cannot change
the request, and she cannot use the signature again. In this build a request
is no more than a right, such as \texttt{write shop/checkout.ts}. A
request that carried data to write would need that data covered by the
signature too, or whoever delivered the proof could choose what was
written. Binding a proof to the connection it arrives on would close this
gap, and this build has no connections: its gate is a function call.

The private key is in the holder's memory, and a process that can read
that memory has it. The gate's table has no size limit, so whoever may
request challenges can fill it for 30 seconds at a time. Two gates that
shared their challenges would need a table with an atomic test-and-set
(Exercise~\ref{ex:anchor-holder-two-gates}). The recorded runs derive their
keys from names and their challenges from a counter, so that each run
repeats byte for byte. Neither is secret, and a gate that protects
anything generates its own at random. Nothing here revokes a card before it
expires (\S\ref{sec:anchor-build-taken-back}), and \pdkernelname{} as a
running service still does not require a card.
\end{pdboundary}
